% =============================================================================
% CHAPTER EXTENSIONS: Ge-on-Si PIN Photodetector
% Islam I. Abdulaal, 2026
%
% SOURCE THESIS ANALYSED: M. G. C. Alasio, "Ge-on-Si photodetectors for
%   silicon photonics: multiphysics modeling and design," Doctoral Dissertation,
%   Politecnico di Torino, June 2024. [cited throughout as \cite{Alasio2024}]
%
% INSTRUCTIONS FOR INTEGRATION:
%   Each block below is labelled with the section of the existing chapter
%   where it should be inserted, and a precise anchor sentence (the last
%   sentence of the preceding paragraph). Paste each block immediately after
%   that anchor, before the next subsection heading. Do NOT alter any
%   existing text. Add \cite{Alasio2024} to the bibliography if not already
%   present (the entry is already in the provided .bib file).
% =============================================================================


% -----------------------------------------------------------------------------
% EXTENSION 1
% Target section : \subsection{Electromagnetic Propagation in an Absorbing
%                  Dispersive Medium} (\cref{subsec:pd_maxwell})
% Insert after   : "...emerges as a corollary of the plane-wave solution to
%                  \cref{eq:helmholtz} rather than an independent empirical
%                  postulate."  (i.e. after \cref{eq:beer_lambert})
% -----------------------------------------------------------------------------

The optical generation rate $G_{\mathrm{opt}}(\mathbf{r})$ defined in
\cref{eq:gen_rate} is equivalently obtained, without recourse to the scalar
Beer-Lambert approximation, from the time-averaged divergence of the Poynting
vector. In a lossless surrounding medium the divergence theorem applied to the
time-averaged Poynting flux $\langle\mathbf{S}\rangle = \tfrac{1}{2}
\mathrm{Re}[\tilde{\mathbf{E}}\times\tilde{\mathbf{H}}^{*}]$ over any closed
surface $\partial\Omega$ enclosing a volume $\Omega$ of absorbing material
gives:
\begin{equation}
    \int_\Omega G_{\mathrm{opt}}(\mathbf{r})\,\hbar\omega\,\mathrm{d}V
    = -\oint_{\partial\Omega}
      \langle\mathbf{S}\rangle\cdot\hat{n}\,\mathrm{d}A,
    \label{eq:gopt_poynting}
\end{equation}
where $\hat{n}$ is the outward unit normal on $\partial\Omega$. The volumetric
generation rate is therefore the local Poynting-flux sink, computed at each
point as:
\begin{equation}
    G_{\mathrm{opt}}(\mathbf{r})
    = -\frac{\nabla\cdot\langle\mathbf{S}(\mathbf{r})\rangle}{\hbar\omega}
    = \frac{\omega\varepsilon_0\,\mathrm{Im}[\tilde{\varepsilon}(\mathbf{r})]
           |\tilde{\mathbf{E}}(\mathbf{r})|^2}{2\,\hbar\omega},
    \label{eq:gopt_poynting_local}
\end{equation}
which reduces identically to \cref{eq:gen_rate} under the plane-wave
approximation $|\tilde{\mathbf{E}}|^2 \propto \exp(-\alpha z)$. The Poynting
formulation is, however, the operationally correct definition when the
three-dimensional electromagnetic field contains interference terms arising
from modal coupling, metal reflections, or end-facet diffraction—all of which
are present in the adiabatic taper region of the device studied here. The FDTD
solver employed in \cref{subsec:pd_fdtd} evaluates $G_{\mathrm{opt}}$ from
\cref{eq:gopt_poynting_local} rather than from the scalar Beer-Lambert
extrapolation; the scalar law serves only as the analytical validation baseline
for the axially integrated profile presented in \cref{fig:gen_rate_map}.


% -----------------------------------------------------------------------------
% EXTENSION 2
% Target section : \subsection{Band Structure of Germanium and the Direct-Gap
%                  Absorption Mechanism} (\cref{subsec:pd_bandstructure})
% Insert after   : "...cannot be used as the absorber. This single
%                  material-physics argument is the entire justification for
%                  the Ge-on-Si heterostructure that forms the subject of this
%                  chapter."
% -----------------------------------------------------------------------------

The electrical consequence of the Ge/Si material interface deserves explicit
treatment because it directly conditions the boundary conditions of the
drift-diffusion problem. Owing to the lattice mismatch ($\Delta a/a \approx
4.2\%$) and the difference in band alignment, the Ge/Si heterojunction
sustains a valence-band discontinuity of $\Delta E_v \approx \SI{0.54}
{\electronvolt}$ and a conduction-band offset $\Delta E_c \approx \SI{-0.14}
{\electronvolt}$~\cite{Alasio2024,Basu2003}. The sign of $\Delta E_c$ is such
that the conduction band of Ge lies \emph{below} that of silicon, producing a
straddling (Type-I like) alignment at the $\Gamma$-valley but a broken-gap
alignment at the $L$-valley due to the indirect-gap energetics. For a
$p^{++}$-Si anode contacting the intrinsic Ge absorber, the large $\Delta E_v$
acts as a potential barrier suppressing hole injection from silicon into
germanium and thereby reducing the minority-carrier diffusion contribution to
the dark current quantified by \cref{eq:jdark_diff}. This barrier is
structurally absent for electrons, which can be injected from the $n^{++}$-Ge
cathode without encountering a conduction-band discontinuity into the Si
waveguide layer below—a fact that motivates the $n$-on-$p$ polarity convention
of the present device. A complication arises, however, in numerical simulation:
the abrupt valence-band step at the Ge/Si interface introduces a singularity in
the quasi-Fermi level that prevents convergence of the self-consistent
drift-diffusion solver unless a thin compositionally graded transition layer is
introduced. Following the approach demonstrated by Van de Walle and
Martin~\cite{Alasio2024} and validated experimentally in the multiphysics
framework of~\cite{Alasio2024}, a linearly interpolated mole-fraction profile
of thickness $\delta_{\mathrm{grad}} \approx \SI{5}{\nano\metre}$ at the
heterojunction removes the discontinuity without perceptibly modifying the
electrostatic potential profile in the intrinsic germanium bulk, and is
accordingly incorporated in the CHARGE simulation of \cref{subsec:pd_charge}.


% -----------------------------------------------------------------------------
% EXTENSION 3
% Target section : \subsection{Carrier Generation and Recombination}
%                  (\cref{subsec:pd_generation})
% Insert after   : "...Radiative band-to-band recombination contributes
%                  negligibly at the doping levels and carrier densities
%                  considered here, because germanium is an indirect-gap
%                  material..."
% -----------------------------------------------------------------------------

The explicit model adopted for radiative recombination is:
\begin{equation}
    U_{\mathrm{rad}} = r_{\mathrm{opt}}\!\left(np - n_i^2\right),
    \label{eq:uradiative}
\end{equation}
where $r_{\mathrm{opt}}$ is the bimolecular optical recombination coefficient.
For germanium at \SI{300}{\kelvin}, $r_{\mathrm{opt}}^{\mathrm{Ge}} \approx
\SI{6.4e-14}{\centi\metre\cubed\per\second}$, approximately four orders of
magnitude smaller than for a direct-gap III-V of comparable bandgap. In the
fully depleted intrinsic region, where $np \approx n_i^2$, the radiative net
rate \cref{eq:uradiative} vanishes identically and \cref{eq:srh_dep}
constitutes the sole generation mechanism in that regime. It is therefore
justified to omit \cref{eq:uradiative} from the analytical dark-current model
of \cref{eq:idark_total} while retaining it in the CHARGE numerical solver as
a consistency check on the total generation-recombination budget. The
simulated generation and recombination rates averaged over the germanium
volume, computed at several reverse-bias values, confirm that SRH exceeds
Auger by more than two orders of magnitude and radiative recombination by more
than five orders of magnitude across the entire operating range $0 \leq
|V_{\mathrm{bias}}| \leq \SI{3}{\volt}$~\cite{Alasio2024}. This hierarchy
justifies the calibration strategy of \cref{subsec:pd_charge} in which the
single free parameter $\tau_0$ is adjusted to match $I_d$ at \SI{-1}{\volt}
without needing independent fits of $r_{\mathrm{opt}}$ or the Auger
coefficients.


% -----------------------------------------------------------------------------
% EXTENSION 4
% Target section : \subsection{Drift-Diffusion Transport: Transient and
%                  Steady-State Formulations} (\cref{subsec:pd_transport})
% Insert after   : "The complete self-consistent system---\cref{eq:poisson}
%                  through \cref{eq:cont_p}---is the foundational PDE set solved
%                  numerically in both the commercial finite-element solver and
%                  the Physics-Informed Neural Network framework described in
%                  \cref{sec:pd_design}."
% -----------------------------------------------------------------------------

\paragraph{Analytical photocurrent transfer function.}
The frequency-domain photocurrent, which is the Fourier-domain counterpart of
the transient continuity equations \cref{eq:cont_n_transient,eq:cont_p_transient},
admits a closed-form expression when diffusion is neglected and the electric
field is assumed uniform at its saturation value throughout the intrinsic
region. Under these conditions, with the harmonic optical stimulus $G_{\mathrm{opt}}
(\mathbf{r},\omega) = G_0(\mathbf{r})\,P_{\mathrm{opt}}(\omega)$ and the
coordinate $z\in[0,W_i]$ measured from the $p^{++}$-contact, the Fourier
amplitudes of the electron and hole current densities at the device terminals
satisfy~\cite{Alasio2024}:
\begin{align}
    J_t(\omega)
    &= \frac{q G_0}{\alpha W_i}\frac{hc}{q\lambda}
       \eta\bigl(1 - R_{\mathrm{refl}}\bigr)P_{\mathrm{opt}}(\omega)
       \times \nonumber \\
    &\quad
      \left\{
        \frac{e^{-\alpha W_i}-1}{\alpha W_i(\alpha W_i - j\omega\tau_{\mathrm{tr,h}})}
        + e^{-\alpha W_i}
          \frac{e^{j\omega\tau_{\mathrm{tr,h}}}-1}
               {j\omega\tau_{\mathrm{tr,h}}(\alpha W_i - j\omega\tau_{\mathrm{tr,h}})}
      \right. \nonumber \\
    &\quad\quad
      \left.
        + \frac{1 - e^{-\alpha W_i}}{\alpha W_i(j\omega\tau_{\mathrm{tr,e}} + \alpha W_i)}
        + \frac{1 - e^{j\omega\tau_{\mathrm{tr,e}}}}{j\omega\tau_{\mathrm{tr,e}}
          (j\omega\tau_{\mathrm{tr,e}} + \alpha W_i)}
      \right\},
    \label{eq:jt_analytical}
\end{align}
where $\tau_{\mathrm{tr,e}} = W_i/v_{\mathrm{d,sat},e}$ and $\tau_{\mathrm{tr,h}}
= W_i/v_{\mathrm{d,sat},h}$ are the electron and hole transit times, and
$R_{\mathrm{refl}}$ is the optical reflectance at the Si/Ge entry facet. The
four terms in \cref{eq:jt_analytical} correspond to holes sweeping toward the
$p^{++}$ contact and electrons toward the $n^{++}$ cathode, each split between
the fraction generated in the first $[\alpha W_i]^{-1}$ of the absorber and the
remaining exponential tail. In the limits $\omega\tau_{\mathrm{tr,h}} \ll 1$
and $\alpha W_i \gg 1$, \cref{eq:jt_analytical} reduces to the rectangular-pulse
sinc-function roll-off that yields $f_{\mathrm{tt}} = 0.45/\tau_{\mathrm{tr,h}}$
as stated in \cref{eq:bw_tt}, confirming that the $0.45$ prefactor is a
consequence of the asymmetric electron-hole velocity ratio
$v_{\mathrm{d,sat},e}/v_{\mathrm{d,sat},h} \approx 1.5$ rather than a purely
geometric constant. The full frequency-dependent transfer function
$H(\omega) = J_t(\omega)/J_t(0)$ derived from \cref{eq:jt_analytical}
is the analytical reference against which the CHARGE small-signal simulation
is validated in \cref{fig:freq_response}.

\paragraph{Scharfetter-Gummel discretisation and numerical stability.}
The finite-element discretisation of \cref{eq:Jn,eq:Jp} on a non-uniform mesh
faces a well-known numerical instability when the local Péclet number
$\mathrm{Pe} = q|\mathbf{E}|\Delta x/(2kT)$ exceeds unity: standard centred
differences oscillate and produce unphysical negative carrier densities. The
Scharfetter-Gummel scheme~\cite{Alasio2024}, which approximates the current
density between adjacent mesh nodes $i$ and $i+1$ as:
\begin{equation}
    J_{n,i+1/2}
    = \frac{qD_n}{\Delta x}\!\left[
        B\!\left(\frac{q\Delta\phi}{kT}\right)n_{i+1}
      - B\!\left(-\frac{q\Delta\phi}{kT}\right)n_i
      \right],
    \label{eq:sg_scheme}
\end{equation}
where $B(u) = u/(e^u - 1)$ is the Bernoulli function and $\Delta\phi =
\phi_{i+1} - \phi_i$, resolves this instability by incorporating the exact
one-dimensional drift-diffusion solution as the interpolation kernel. The
Bernoulli function recovers the centred-difference approximation for
$|u| \ll 1$ (low-field, diffusion-dominated) and the upstream-weighted
approximation for $|u| \gg 1$ (high-field, drift-dominated), providing
second-order accuracy and unconditional positivity across the full range of
electric field magnitudes encountered in the \SI{30}{\kilo\volt\per\centi\metre}
regime of the present device. The CHARGE solver applies \cref{eq:sg_scheme}
on each edge of the adaptive tetrahedral mesh, with the mesh refinement
criterion requiring that $\mathrm{Pe} < 3$ everywhere in the germanium bulk;
this criterion is the operative test for mesh convergence rather than an
absolute cell-size constraint, because the local field magnitude varies by more
than an order of magnitude between the ohmic contact tails and the peak of the
depletion-region field.


% -----------------------------------------------------------------------------
% EXTENSION 5
% Target section : \subsubsection{Capacitance and Series Resistance}
%                  (\cref{subsubsec:pd_cres})
% Insert after   : "...from which $f_{\mathrm{RC}} = (2\pi\tau_{\mathrm{RC}})^{-1}$."
% -----------------------------------------------------------------------------

The small-signal impedance of the device provides a direct, bias-resolved
measurement of the individual circuit components that enter $\tau_{\mathrm{RC}}$.
The admittance $Y(\omega) = I_{\mathrm{AC}}/V_{\mathrm{AC}}$ extracted from
the CHARGE small-signal analysis exhibits two distinct plateau levels in
$\mathrm{Im}[Y(\omega)]/(2\pi f)$ as a function of frequency: a low-frequency
plateau at $C_1 + C_2$ and a high-frequency plateau at $C_2$ alone, where
$C_1$ is the intrinsic junction capacitance and $C_2$ is the parasitic
contribution from the first interconnect metal layer. For the vertical
$n$-$i$-$p$ architecture under study, with a junction area
$A_j = W_m \times L = \SI{1}{\micro\metre} \times \SI{6}{\micro\metre}$ and
the reference device parameters of \cref{tab:layer_stack}, the parallel-plate
estimate $C_1 \approx \varepsilon_0 \varepsilon_{\mathrm{Ge}} A_j / W_i
\approx \SI{1.8}{\femto\farad}$ is verified by the low-to-high-frequency
plateau difference extracted numerically. The parasitic contribution $C_2$,
arising from the capacitive coupling between the first copper metal layer and
the underlying silicon platform, is not negligible relative to $C_1$: it adds
a comparable fraction to the total capacitance $C_{\mathrm{tot}} = C_1 + C_2$
that enters $\tau_{\mathrm{RC}}$~\cite{Alasio2024}. The critical implication
is that the effective series resistance $R$ extracted from
$\mathrm{Re}[Y(\omega)]$ is essentially independent of whether or not the
first metal layer is included in the simulation, whereas $C_2$ changes
substantially: removing the copper layer reduces $C_2$ by more than
\SI{35}{\percent}. Electrode layout optimisation must therefore simultaneously
minimise the lateral conduction path $\ell_a$ in \cref{eq:rs_anode} and
reduce the overlap area between the first metal layer and the doped silicon
platform, because both $R_s$ and $C_2$ contribute multiplicatively to
$\tau_{\mathrm{RC}}$.


% -----------------------------------------------------------------------------
% EXTENSION 6
% Target section : \subsection{Optical Simulation: Three-Dimensional FDTD}
%                  (\cref{subsec:pd_fdtd})
% Insert after   : "Perfectly matched layer boundary conditions are applied at
%                  all domain faces to suppress artificial reflections."
% -----------------------------------------------------------------------------

The PML implementation employs sixteen absorbing cells on each domain
boundary, a number determined by convergence of the reflected power to below
$-60\,\mathrm{dB}$ of the incident power in a calibration run with a uniform
germanium slab. The temporal stability of the Yee algorithm requires the
Courant-Friedrichs-Lewy condition $c\Delta t < \Delta x$, which in three
dimensions with isotropic mesh spacing $\Delta$ tightens to the Courant
stability criterion:
\begin{equation}
    \Delta t < \frac{\Delta}{c\sqrt{3}},
    \label{eq:cfl_3d}
\end{equation}
so that the time step is set to $\Delta t = 0.99\,\Delta/(c\sqrt{3})$ for
all simulations. The spatial mesh size is independently constrained by the
sampling requirement $\Delta < \lambda_{\min}/(10\sqrt{\varepsilon_{\max}})$,
where $\lambda_{\min} = \SI{1310}{\nano\metre}$ is the free-space wavelength
and $\varepsilon_{\max} = \varepsilon_{\mathrm{Ge}} = 16.2$ is the maximum
relative permittivity in the domain; this yields $\Delta < \SI{32}{\nano\metre}$
and the \SI{20}{\nano\metre} override mesh in the germanium region satisfies
this criterion with a safety margin of 1.6. The combined optical simulation
(FDTD) and electrical simulation (drift-diffusion, including the small-signal
AC sweep for capacitance extraction) require approximately \SI{11}{\hour} of
wall-clock time on a shared-memory workstation for the VPIN geometry of
\cref{tab:layer_stack}, of which roughly \SI{3}{\hour} are consumed by the
FDTD step and \SI{8}{\hour} by the steady-state and frequency-response
electrical solves~\cite{Alasio2024}. This computational budget constrains the
number of independent parameter sweeps that can be performed within the scope
of this thesis to the absorber length sweep of \cref{subsec:pd_optical_results}
and the bias-sweep of \cref{subsec:pd_elec_profiles}, and motivates the PINN
surrogate discussed in \cref{subsec:pd_pinns} as a strategy for accelerating
future design-space exploration.


% -----------------------------------------------------------------------------
% EXTENSION 7
% Target section : \subsection{Electrical Simulation: Three-Dimensional
%                  Drift-Diffusion} (\cref{subsec:pd_charge})
% Insert after   : "Ohmic boundary conditions are imposed at the cathode
%                  (biased to the desired reverse voltage $V < 0$) and anode
%                  contacts (grounded, $V = 0$)."
% -----------------------------------------------------------------------------

The Canali velocity-saturation model adopted in \cref{eq:mobility_sat}
follows the original parametrisation of Canali \textit{et al.}~(1975), in
which the exponent $\beta$ and the prefactor $\alpha_{\mathrm{C}}$ are
empirical fitting constants that capture the transition from ohmic to
saturation regime as a function of the local field magnitude. Explicitly:
\begin{equation}
    \mu_{n,p}(|\mathbf{E}|)
    = \frac{(\alpha_{\mathrm{C}}+1)\,\mu_{n,p}^{(0)}}
      {\left[\alpha_{\mathrm{C}}
        + \!\left(1 + \!\left(\frac{(\alpha_{\mathrm{C}}+1)\,
          \mu_{n,p}^{(0)}\,|\mathbf{E}|}{v_{\mathrm{d,sat}}}\right)^{\!\beta}
        \right)^{1/\beta}\right]},
    \label{eq:canali_full}
\end{equation}
which reduces to \cref{eq:mobility_sat} with $\beta = 2$ and $\alpha_{\mathrm{C}}
= 0$~\cite{Alasio2024}. The low-field mobilities of germanium adopted in the
present simulation are $\mu_n^{(0)} = \SI{3900}{\centi\metre\squared\per\volt
\per\second}$ and $\mu_p^{(0)} = \SI{1900}{\centi\metre\squared\per\volt
\per\second}$, corresponding to the bulk intrinsic values at \SI{300}{\kelvin};
the saturation velocities $v_{\mathrm{d,sat},e} = \SI{6e6}{\centi\metre\per
\second}$ and $v_{\mathrm{d,sat},h} = \SI{4e6}{\centi\metre\per\second}$
enter \cref{eq:canali_full} as the asymptotic limits. The low-mobility carrier,
holes, reaches $v_{\mathrm{d,sat},h}$ at a field of approximately
$E_{\mathrm{sat}} = v_{\mathrm{d,sat},h}/\mu_p^{(0)} \approx \SI{21}{\kilo
\volt\per\centi\metre}$ in the linear approximation; the electric field profile
in \cref{fig:efield} confirms that the intrinsic layer field exceeds
$E_{\mathrm{sat}}$ by more than a factor of 1.4 across its full thickness at
\SI{-1}{\volt}, validating velocity saturation throughout the active volume.
An important practical consequence of this confirmation is that the
drift-diffusion model operates in a regime where $\mu_p(|\mathbf{E}|)$ has
already converged to $v_{\mathrm{d,sat},h}/|\mathbf{E}|$ and is therefore
insensitive to the exact value of $\mu_p^{(0)}$: a $\pm 20\%$ variation in
$\mu_p^{(0)}$ changes $f_{\mathrm{tt}}$ by less than \SI{1}{\percent},
confirming the robustness of the transit-time estimate in \cref{eq:bw_tt}.


% -----------------------------------------------------------------------------
% EXTENSION 8
% Target section : \section{Discussion} (\cref{sec:pd_discussion})
% Insert as a new \paragraph immediately after the existing paragraph titled
%   "Bandwidth and the intrinsic-layer thickness trade-off."
% -----------------------------------------------------------------------------

\paragraph{Electric field screening under strong optical injection.}
The analysis of \cref{sec:pd_results} is conducted at an input optical power
of \SI{100}{\micro\watt}, corresponding to the low-injection regime in which
the density of photogenerated carriers is small relative to the background
doping and the electrostatic solution is insensitive to illumination level.
A physically distinct and operationally significant regime emerges at higher
optical powers, where the photogenerated carrier density within the intrinsic
germanium layer becomes comparable to or exceeds the background $n_i$ and
modifies the local electric field through its own space charge. This
\emph{electric field screening} mechanism was characterised experimentally and
reproduced by the coupled FDTD/drift-diffusion framework of~\cite{Alasio2024},
which showed that increasing the input power from $-2\,\mathrm{dBm}$ to
$+3\,\mathrm{dBm}$ reduces the peak electric field in the intrinsic germanium
by a factor of three at a longitudinal position one micrometre from the taper
exit—a reduction sufficient to push the local field below the carrier saturation
threshold $E_{\mathrm{sat}} \approx \SI{7}{\kilo\volt\per\centi\metre}$ and
thereby restore sub-saturation carrier velocities. The consequence in terms of
device metrics is a monotonic reduction of $f_{3\,\mathrm{dB}}$ with increasing
optical power of more than \SI{20}{\percent} between the lowest and highest
powers tested. Within the analytical framework of \cref{eq:bw_tt}, this
corresponds to an effective increase of $W_i$ (the carriers transit a longer
effective thickness before reaching their collection electrodes at reduced
velocity), which is structurally equivalent to a reduction of
$v_{\mathrm{d,sat},h}$ in \cref{eq:bw_tt} from the geometric perspective of
the bandwidth model. Quantitatively, for the device geometry of \cref{tab:layer_stack}
at \SI{-1}{\volt}, the screening onset occurs near
$P_{\mathrm{opt}} \approx \SI{200}{\micro\watt}$, consistent with the
condition $G_{\mathrm{opt,peak}}\tau_{\mathrm{SRH}} \sim n_i$, which requires:
\begin{equation}
    P_{\mathrm{screening}}
    \approx \frac{n_i\,\hbar\omega\,A_j\,W_i}{\tau_0\,\alpha(\lambda)\,\eta_c},
    \label{eq:p_screening}
\end{equation}
where $\tau_0 = \SI{5}{\nano\second}$ and $\alpha = \SI{2.1e3}{\per
\centi\metre}$ as calibrated in \cref{subsec:pd_charge}. Substituting the
device parameters gives $P_{\mathrm{screening}} \approx \SI{340}{\micro\watt}$,
in order-of-magnitude agreement with the numerical observation and confirming
that the present operating point of \SI{100}{\micro\watt} is safely below the
onset. The screening effect is not captured by the analytical model of
\cref{eq:bw_tt} because that model assumes the electric field is independent of
the optical stimulus; it is, however, fully accounted for by the self-consistent
CHARGE simulation, and \cref{eq:p_screening} provides an explicit design
inequality for the maximum usable optical power at a given junction geometry
and reverse bias.

\paragraph{Doping implantation geometry and its effect on field uniformity.}
A further geometric design variable not addressed in the primary metric
framework of \cref{sec:pd_concept} is the lateral extent of the highly doped
cathode region within the germanium mesa, parameterised as
$W_{\mathrm{doping}}$ in~\cite{Alasio2024}. When $W_{\mathrm{doping}} <
W_m$, the electric field beneath the undoped corners of the germanium surface is
reduced relative to the centre of the absorber, because those peripheral regions
do not benefit from the vertical field enforced by the overlying metal contact.
The CHARGE simulations of~\cite{Alasio2024} quantify this effect through a
parametric sweep of $W_{\mathrm{doping}}$ at fixed $W_m$: for the reference
geometry, reducing $W_{\mathrm{doping}}$ from $W_m$ to $0.75\,W_m$ decreases
$f_{3\,\mathrm{dB}}$ by more than \SI{60}{\percent} at \SI{-0.8}{\volt} and
by approximately \SI{30}{\percent} at \SI{-2}{\volt}, with an abrupt
performance cliff near $W_{\mathrm{doping}} \approx 0.75\,W_m$ that coincides
with the onset of sub-saturation carrier velocities in the corner regions.
This result adds a third geometric constraint to the absorber design
optimisation beyond the absorber length $L$ and thickness $W_i$ that govern
\cref{eq:abs_fraction,eq:bw_tt}: the ratio $W_{\mathrm{doping}}/W_m$ must
approach unity to ensure that the electric field is vertically oriented
everywhere within the active volume, a condition that simultaneously maximises
the carrier collection efficiency by suppressing lateral diffusion toward the
lower-field corners. For the U-shaped electrode architecture of the present
device, this requirement is naturally satisfied because the anode metal wraps
the mesa perimeter and the cathode implantation covers the full top surface;
its practical significance is greatest for alternative electrode topologies
in which the metal contact area is constrained by isolation requirements.

\paragraph{Quantitative bandwidth-responsivity trade-off and energy consumption.}
The coupled dependence of $f_{3\,\mathrm{dB}}$ and $\mathcal{R}$ on the
absorber thickness $W_i$ can be quantified through the CHARGE parameter sweep
of~\cite{Alasio2024}, which establishes that reducing $W_i$ from \SI{800}{\nano
\metre} to \SI{400}{\nano\metre} raises $f_{\mathrm{tt}}$ in \cref{eq:bw_tt}
by \SI{25}{\giga\hertz} while simultaneously decreasing $\mathcal{R}$ by
approximately \SI{9}{\percent}. The exchange rate between bandwidth and
responsivity in this regime is therefore approximately
\SI{2.8}{\giga\hertz\per\percent} of responsivity sacrificed, a figure that
constitutes a device-specific constant of the vertical n-i-p architecture at
the O-band. An additional metric not previously discussed, the energy consumed
per detected bit, $e_c\,[\mathrm{J/bit}]$, connects the device parameters to
the power budget of the photonic integrated circuit. Under the assumption that
the photocurrent is insensitive to the applied bias (satisfied for all
reverse-biased conditions within the present operating range), $e_c$ is
proportional to the product of the bias voltage and the photocurrent per unit
optical power, hence $e_c \propto |V_{\mathrm{bias}}|\,\mathcal{R}$. For the
reference geometry at a \SI{10}{\giga\baud} symbol rate and \SI{-2}{\volt}
bias, the estimated energy per bit is $e_c \approx \SI{35}{\femto\joule\per
\bit}$~\cite{Alasio2024}, a value competitive with the state-of-the-art for
integrated germanium photodetectors and consistent with the power consumption
requirements of next-generation silicon photonic data-centre interconnects.
The connection between device-level $e_c$ and the system-level link power
budget complements the shot-noise-limited sensitivity of \cref{eq:sens_result}
as a second design axis: whereas sensitivity bounds the minimum detectable
power, $e_c$ bounds the maximum permissible photocurrent for a given power
envelope, and the two constraints together define the feasible operating window
for the receiver.

\paragraph{Lateral junction scaling limit and carrier velocity overshoot.}
The foregoing discussion has been confined to the vertical junction topology.
The bandwidth ceiling of \SI{265}{\giga\hertz} achieved by the lateral
FinFET-like device of Lischke~\textit{et al.}~\cite{Lischke2021} introduces
a physically distinct regime that illuminates, by contrast, the mechanisms
operating in the vertical architecture. In the lateral geometry, the intrinsic
germanium width $W_{\mathrm{Ge}}$ plays the role of $W_i$ in \cref{eq:bw_tt},
but can be scaled to $W_{\mathrm{Ge}} \approx \SI{100}{\nano\metre}$ without
modifying the optical absorption path length $L$. At this scale, the
drift-diffusion saturation-velocity model that underpins \cref{eq:bw_tt}
becomes quantitatively inadequate: full-band Monte Carlo simulations of the
lateral device show that holes overshoot their drift-diffusion saturation
velocity by more than \SI{40}{\percent} for approximately the first
\SI{50}{\nano\metre} of their trajectory—representing half the absorber
width—before relaxing back to $v_{\mathrm{d,sat},h}$ as phonon scattering
equilibrates the hot-carrier distribution~\cite{Alasio2024}. The effect of
this transient overshoot is an effective transit time shorter than
$W_{\mathrm{Ge}}/v_{\mathrm{d,sat},h}$ would predict, partially explaining
the experimentally observed bandwidth exceeding the drift-diffusion transit-time
ceiling. A second anomaly specific to the aggressively scaled lateral device
is the emergence of impact ionisation: as $W_{\mathrm{Ge}}$ decreases below
\SI{100}{\nano\metre}, the ratio of the avalanche generation rate to the total
optical generation rate approaches and may exceed unity, meaning that
impact-ionisation-generated carriers contribute a current component comparable
in magnitude to the photocurrent itself~\cite{Alasio2024}. This regime,
characterised by avalanche multiplication of the photogenerated seed carriers
within a sub-micrometre transit length, lies entirely outside the validity of
the drift-diffusion model and requires a full-band Monte Carlo treatment for
quantitative prediction. For the vertical architecture of the present chapter,
where $W_i = \SI{350}{\nano\metre}$ and the saturation field is exceeded
throughout the depletion region at \SI{-1}{\volt}, neither overshoot nor
avalanche multiplication is significant at the operating point, but both
effects become relevant if $W_i$ is further reduced below \SI{200}{\nano
\metre} in pursuit of additional transit-time bandwidth—providing an explicit
physical lower bound on the absorber thickness for which the drift-diffusion
framework of \cref{subsec:pd_transport} remains the appropriate transport model.


% -----------------------------------------------------------------------------
% EXTENSION 9
% Target section : \section{Device-to-System Projection} (\cref{sec:pd_system})
% Insert as a new \subsection immediately before \cref{subsec:pd_sens},
%   to be entitled:
%   \subsection{Doping Configuration and Its Impact on the Receiver Dark Current
%               Floor}
% -----------------------------------------------------------------------------

\subsection{Doping Configuration and Its Impact on the Receiver Dark Current
            Floor}
\label{subsec:pd_polarity}

The dark current $I_d$ that enters the shot-noise expression \cref{eq:shot}
and the sensitivity formula \cref{eq:sens_full} is not uniquely determined by
the junction geometry: for fixed $W_i$, $A_j$, and $\tau_0$, the doping
polarity of the heterostructure—whether the germanium is made $n$-type and the
silicon platform $p$-type ($n$-on-$p$), or vice versa ($p$-on-$n$)—introduces
an asymmetry in the dark current through the heterojunction band alignment at
the Ge/Si interface. The $n$-on-$p$ configuration places the graded Ge/Si
transition within the Ge depletion region, where the valence-band discontinuity
$\Delta E_v$ is traversed by thermally generated holes. The associated
potential notch modifies the effective barrier for minority-carrier generation
at the interface, raising $J_{\mathrm{surf}}$ in \cref{eq:jdark_surf} above
the value predicted by the ideal planar junction model~\cite{Alasio2024}. In
the $p$-on-$n$ configuration, where the silicon platform is $n$-doped and the
germanium cathode is $p$-doped, the same $\Delta E_v$ acts as an additive
barrier to interface hole generation, and the resulting dark current is
correspondingly lower. The practical consequence for the system projection of
\cref{subsec:pd_sens} is that for a fixed reverse bias of \SI{-1}{\volt}, the
$p$-on-$n$ device exhibits an $I_d$ that is a factor of approximately two to
three lower than the $n$-on-$p$ counterpart, a difference that translates
through \cref{eq:sens_dark} into a sensitivity improvement of approximately
\SI{1.5}{\decibel} in the dark-current-limited regime. The present device
adopts the $p^{++}$-Si/$i$-Ge/$n^{++}$-Ge ($p$-on-$n$ with respect to the
Si platform) polarity precisely because it minimises the interface generation
contribution to $I_d$ while maintaining the CMOS process compatibility of a
$p$-doped silicon-on-insulator layer. This polarity choice is therefore not an
arbitrary convention but a deliberate engineering decision grounded in the
heterojunction band alignment physics of \cref{subsec:pd_bandstructure}.


% -----------------------------------------------------------------------------
% EXTENSION 10
% Target section : \section{Conclusion} (\cref{sec:pd_conclusion})
% Insert as an additional paragraph at the end of the Conclusion, before the
%   sentence "Future experimental work should address..."
% -----------------------------------------------------------------------------

The multiphysics methodology demonstrated here—coupling a three-dimensional
FDTD optical solver to a self-consistent drift-diffusion electrical solver and
cross-validating with a mesh-free PINN—defines a simulation chain that is
more stringent than any of the individual solvers alone. The chain exposes an
important hierarchy of physical mechanisms whose relative importance shifts as
the device geometry is scaled: in the \SI{350}{\nano\metre}-thick vertical
absorber studied here, the SRH dark current and the saturation-velocity
transit time are the dominant performance limiters, and both are well-described
by the continuum drift-diffusion formalism. As $W_i$ is reduced below
\SI{200}{\nano\metre}, however, two physically distinct effects invalidate the
continuum description: transient carrier velocity overshoot elevates the
effective transit velocity above $v_{\mathrm{d,sat},h}$, while impact
ionisation introduces a positive feedback between photocurrent and carrier
multiplication that creates a non-linear photocurrent-power relationship
outside the responsivity model of \cref{eq:responsivity}. A complete design
framework for sub-200-nm absorbers therefore requires a full-band Monte Carlo
transport solver as a complement to the drift-diffusion chain employed
here~\cite{Alasio2024}. Additionally, the electric field screening effect of
\cref{sec:pd_discussion} establishes an explicit optical power ceiling
$P_{\mathrm{screening}}$ from \cref{eq:p_screening} beyond which the
bandwidth model of \cref{eq:bw_total} loses quantitative validity; this
ceiling must be incorporated as a hard constraint in any link power budget
that specifies the maximum launch power at the chip input. The physical
intuitions and analytical inequalities developed in this chapter—covering
\cref{eq:alpha_kappa} through \cref{eq:p_screening}—constitute a
self-consistent toolkit for navigating these trade-offs in the design of the
next generation of Ge-on-Si receivers.


% =============================================================================
% SUGGESTED ADDITIONAL FIGURES AND TABLES
% (place within the corresponding sections as indicated)
% =============================================================================

% [Suggested Figure -- Section 4.2 (Device in dark) / \cref{sec:pd_results}]
% Description : Side-by-side band diagrams at equilibrium for the p-on-n and
%   n-on-p configurations, plotted along a vertical 1D cut from the top metal
%   contact through the n++-Ge cathode, i-Ge absorber, and p++-Si platform.
%   Both should show Ec(z), Ev(z), and EF, with the Ge/Si valence-band
%   discontinuity Delta_Ev annotated and the graded-interface region shaded.
% Purpose : Demonstrates how the polarity choice affects the interface barrier
%   height for hole generation and motivates the p-on-n configuration adopted
%   in the present device through the dark-current argument of
%   \cref{subsec:pd_polarity}. Connects to the discussion in Extension 9.

% [Suggested Figure -- Section 4.3 (Device under illumination) / \cref{subsec:pd_bw_results}]
% Description : Normalised electron drift velocity |v_n(x,y,z)|/v_sat in the
%   transverse (xy) cross-section of the Ge absorber at z = L/4 from the
%   taper exit, plotted for two reverse-bias values: -1 V (present operating
%   point) and 0 V (zero-bias). The saturation contour (|v|/v_sat = 1) should
%   be overlaid as a dashed line. A companion panel should show the equivalent
%   hole velocity map.
% Purpose : Provides direct visual evidence that velocity saturation is
%   achieved across the entire absorber cross-section at -1 V (validating
%   eq:bw_tt) while showing the residual sub-saturation pockets at zero bias
%   relevant to the zero-bias bandwidth of \cref{sec:pd_discussion}.

% [Suggested Table -- Section 4.5 (Design guidelines) / \cref{sec:pd_discussion}]
% Description : A table with columns: W_i (nm) | f_tt (GHz) | C_j (fF) |
%   f_RC (GHz) | f_3dB (GHz) | R (A/W) | e_c (fJ/bit), sweeping W_i from
%   200 nm to 600 nm in 100 nm steps at fixed A_j = 6 um^2 and V = -1 V.
%   Values for R and e_c should be taken from or extrapolated consistently
%   with the Beer-Lambert model of eq:qe_chain. Mark the present design point
%   (W_i = 350 nm) with a bold row.
% Purpose : Consolidates the bandwidth-responsivity-energy trade-off of
%   Extension 8 into a single reference table, replacing the qualitative
%   narrative with specific numerical predictions that can be directly
%   compared with future experimental measurements.

% [Suggested Figure -- Section 5 / \cref{sec:pd_discussion}]
% Description : Electro-optic frequency response |H(f)|^2 (dB) from the
%   analytical transfer function of eq:jt_analytical, plotted for three bias
%   conditions: V = 0 V, -0.5 V, -1 V. Each curve should be labelled with the
%   effective hole transit velocity v_h/v_sat (ratio to saturation), derived
%   from the Canali model at the mean electric field for that bias. A dashed
%   line at -3 dB marks the bandwidth for each condition.
% Purpose : Explicitly connects the bias dependence of f_3dB to the carrier
%   velocity ratio, making the saturation-velocity assumption of eq:bw_tt
%   visually testable and bounding the error in the analytical model for
%   below-saturation conditions relevant to zero-bias operation.

% =============================================================================
% NEW BIBLIOGRAPHY ENTRY (already in the provided .bib file as \cite{Alasio2024})
% No action required; included here for completeness.
% @phdthesis{Alasio2024,
%   author  = {Alasio, Matteo Giovanni Carmelo},
%   title   = {Ge-on-Si photodetectors for silicon photonics: multiphysics
%              modeling and design},
%   school  = {Politecnico di Torino},
%   type    = {Doctoral Dissertation},
%   address = {Torino, Italy},
%   month   = jun,
%   year    = {2024},
%   url     = {https://hdl.handle.net/11583/2990835}
% }
% =============================================================================
